% ============================================================
%  Chapter 17 — Diagonalization and Eigenvalues
% ============================================================
\section{Diagonalization and Eigenvalues}
\label{sec:diagonalization}

A fundamental question in linear algebra is whether a linear operator can be
represented, in some basis, by a diagonal matrix.  This chapter develops the
theory of eigenvalues and eigenspaces, which provides the complete answer.

% ---- Aim of diagonalization ---------------------------------
\begin{tipbox}[Remark --- Aim of Diagonalization]
\label{rem:aim-diag}
Given an $n$-dimensional vector space $V$ over $\F$ and a linear operator
$f:V\to V$, the goal is to find an ordered basis $\mathcal{B}$ for $V$ such
that
\[
    [f]_{\mathcal{B}} = \operatorname{diag}(\lambda_1,\lambda_2,\ldots,\lambda_n).
\]
\end{tipbox}

% ---- Diagonalizable l.t. ------------------------------------
\begin{defbox}[Definition --- Diagonalizable Linear Transformation]
\label{def:diagonalizable}
A linear transformation $f:V\to V$ is called \textbf{diagonalizable} if
there exists an ordered basis $\mathcal{B}$ for $V$ such that
\[
    [f]_{\mathcal{B}} = \operatorname{diag}(\lambda_1,\ldots,\lambda_n)
    = \begin{bmatrix}
        \lambda_1 & 0 & \cdots & 0 \\
        0 & \lambda_2 & \cdots & 0 \\
        \vdots & & \ddots & \vdots \\
        0 & 0 & \cdots & \lambda_n
      \end{bmatrix}.
\]
A matrix $A\in\F^{n\times n}$ is diagonalizable if there exists an
invertible $P\in\F^{n\times n}$ such that $P^{-1}AP$ is diagonal, i.e.\
$A=PDP^{-1}$ for some diagonal matrix $D$.
\end{defbox}

% ---- Eigenvector --------------------------------------------
\begin{defbox}[Definition --- Eigenvector]
\label{def:eigenvector}
A \textbf{nonzero} vector $v\in V$ is called an \textbf{eigenvector} of $f$
if there exists a scalar $\lambda\in\F$ such that
\[
    f(v) = \lambda v.
\]
\end{defbox}

% ---- Eigenvalue and eigenspace ------------------------------
\begin{defbox}[Definition --- Eigenvalue and Eigenspace]
\label{def:eigenvalue-eigenspace}
A scalar $\lambda\in\F$ is called an \textbf{eigenvalue} of $f$ if
\[
    \Ker(f - \lambda\,\mathrm{id}) \neq \{0\},
\]
equivalently, if there exists $v\in V\setminus\{0\}$ such that
$f(v)=\lambda v$.

The set
\[
    V_\lambda = \{v\in V \mid f(v)=\lambda v\}
\]
is called the \textbf{eigenspace} of $f$ associated with $\lambda$
(or \emph{relative to} $\lambda$).  Its nonzero elements are the
eigenvectors corresponding to $\lambda$.
\end{defbox}

\begin{tipbox}[Remark --- Eigenspace as a Kernel]
For any scalar $\lambda\in\F$ and linear operator $f:V\to V$:
\[
    f(v)=\lambda v
    \;\Longleftrightarrow\;
    (f-\lambda\,\mathrm{id})(v)=0
    \;\Longleftrightarrow\;
    v\in\Ker(f-\lambda\,\mathrm{id}).
\]
Hence $V_\lambda = \Ker(f-\lambda\,\mathrm{id})$, which is a
\textbf{subspace} of $V$.
\end{tipbox}

% ---- Equivalent conditions for eigenvalue -------------------
\begin{thmbox}[Theorem --- Equivalent Conditions for Eigenvalues]
\label{thm:eigenvalue-conditions}
The following are equivalent:
\begin{enumerate}[label=\textbf{(\roman*)}]
    \item $\lambda$ is an eigenvalue of $f$.
    \item $f-\lambda\,\mathrm{id}$ is not a monomorphism (i.e.\ not injective).
    \item $\det(f-\lambda\,\mathrm{id})=0$.
\end{enumerate}
\end{thmbox}

% ---- Characteristic polynomial ------------------------------
\subsection{The Characteristic Polynomial}

\begin{defbox}[Definition --- Characteristic Polynomial]
\label{def:char-poly}
Let $f:V\to V$ be a linear operator on a vector space $V$ over $\F$, and
let $\mathcal{B}$ be an ordered basis for $V$.  The \textbf{characteristic
polynomial} of $f$ is the monic polynomial
\[
    \chi_f(x) = \det\!\bigl(x\cdot\mathrm{id} - [f]_{\mathcal{B}}\bigr).
\]
\end{defbox}

\begin{warnbox}[Notation Notice --- Convention for the Characteristic Polynomial]
The notes define the characteristic polynomial as
$\chi_f(x)=\det(x\cdot\mathrm{id}-[f]_{\mathcal{B}})$.
Some textbooks use the opposite sign convention
$\chi_f(x)=\det([f]_{\mathcal{B}}-x\cdot\mathrm{id})$, which gives the
same roots but may differ in sign of coefficients.  \textbf{These notes
follow the convention $\chi_f(x)=\det(x\cdot I - A)$.}
\end{warnbox}

\begin{corbox}[Corollary --- Eigenvalues are Roots of the Characteristic Polynomial]
\label{cor:eigenvalues-roots}
Let $f:V\to V$ be a linear operator on a finite-dimensional vector space
$V$ over $\F$.  A scalar $\lambda\in\F$ is an eigenvalue of $f$ if and
only if $\lambda$ is a root of $\chi_f(x)$.
\end{corbox}

\begin{propbox}[Proposition --- Characteristic Polynomial is Basis-Independent]
\label{prop:chi-basis-independent}
$\chi_f(x)$ does not depend on the choice of ordered basis $\mathcal{B}$.
\end{propbox}

\begin{proofbox}[Proof --- Characteristic Polynomial is Basis-Independent]
If $\mathcal{B}$ and $\mathcal{B}'$ are two ordered bases, then
$[f]_{\mathcal{B}'}=C^{-1}[f]_{\mathcal{B}}C$ for some invertible $C$.
Hence
\[
    \det(xI-[f]_{\mathcal{B}'})
    = \det(C^{-1}(xI-[f]_{\mathcal{B}})C)
    = \det(xI-[f]_{\mathcal{B}}).
    \hfill\blacksquare
\]
\end{proofbox}

% ---- Worked example -----------------------------------------
\begin{exbox}[Example --- Characteristic Polynomial, Eigenvalues, and Eigenspaces]
\label{ex:char-poly}
Let $\mathcal{B}=\{v_1,v_2\}$ be an ordered basis for $\R^2$ and define
$f:\R^2\to\R^2$ by $f(v_1)=2v_1+2v_2$ and $f(v_2)=3v_1+v_2$.

\medskip\noindent
\textbf{Matrix:}
$[f]_{\mathcal{B}} = \begin{pmatrix}2&3\\2&1\end{pmatrix}$.

\medskip\noindent
\textbf{Characteristic polynomial:}
\[
    \chi_f(x)
    = \det\!\begin{pmatrix}x-2&-3\\-2&x-1\end{pmatrix}
    = (x-2)(x-1)-6
    = x^2-3x-4
    = (x-4)(x+1).
\]

\medskip\noindent
\textbf{Eigenvalues:} $\lambda_1=4$, $\lambda_2=-1$.

\medskip\noindent
\textbf{Eigenspace $V_4$:}
\[
    V_4 = \Ker\!\begin{pmatrix}-2&3\\2&-3\end{pmatrix}
        = \Ker\!\begin{pmatrix}-2&3\\0&0\end{pmatrix}
        = \operatorname{span}\!\{(3,2)\}
        \;\Longrightarrow\;
        V_4 = \operatorname{span}\{3v_1+2v_2\}.
\]

\medskip\noindent
\textbf{Eigenspace $V_{-1}$:}
$V_{-1}=\operatorname{span}\{v_1-v_2\}$ (similar computation).

\medskip\noindent
\textbf{Conclusion:}
$\R^2 = V_4\oplus V_{-1}$, and the basis
$\mathcal{B}'=\{3v_1+2v_2,\;v_1-v_2\}$ satisfies
\[
    [f]_{\mathcal{B}'} = \begin{pmatrix}4&0\\0&-1\end{pmatrix},
    \qquad
    [f]_{\mathcal{B}} = C(\mathcal{B}',\mathcal{B})\,[f]_{\mathcal{B}'}\,C(\mathcal{B},\mathcal{B}'),
\]
where $C(\mathcal{B}',\mathcal{B})=\begin{pmatrix}3&1\\2&-1\end{pmatrix}$
and $C(\mathcal{B},\mathcal{B}')=C(\mathcal{B}',\mathcal{B})^{-1}$.
\end{exbox}

% ---- Intersection of eigenspaces ----------------------------
\begin{propbox}[Proposition --- Eigenspaces for Distinct Eigenvalues Intersect Trivially]
\label{prop:eigenspace-intersection}
Let $f:V\to V$ be a linear operator on a finite-dimensional vector space
$V$ over $\F$.  If $\lambda_1\neq\lambda_2$ are distinct eigenvalues of
$f$, then
\[
    V_{\lambda_1}\cap V_{\lambda_2} = \{0\}.
\]
\end{propbox}

\begin{corbox}[Corollary --- Sum of Two Eigenspaces is Direct]
Under the same hypotheses, $V_{\lambda_1}+V_{\lambda_2}=V_{\lambda_1}\oplus V_{\lambda_2}$.
\end{corbox}

\begin{thmbox}[Theorem --- Sum of Eigenspaces is Direct]
\label{thm:eigenspaces-direct}
Let $\lambda_1,\lambda_2,\ldots,\lambda_k$ be \emph{distinct} eigenvalues
of $f:V\to V$.  Then
\[
    V_{\lambda_1}+V_{\lambda_2}+\cdots+V_{\lambda_k}
    = V_{\lambda_1}\oplus V_{\lambda_2}\oplus\cdots\oplus V_{\lambda_k}.
\]
\end{thmbox}

% ---- Multiplicities -----------------------------------------
\begin{defbox}[Definition --- Geometric and Algebraic Multiplicity]
\label{def:multiplicities}
Let $\lambda$ be an eigenvalue of $f$.
\begin{enumerate}[label=\textbf{(\alph*)}]
    \item The \textbf{geometric multiplicity} of $\lambda$ is
          $\dime V_\lambda = \dime\Ker(f-\lambda\,\mathrm{id})$.
    \item The \textbf{algebraic multiplicity} of $\lambda$ is
          $\operatorname{mult}(\lambda,\chi_f(x))$, the multiplicity
          of $\lambda$ as a root of the characteristic polynomial.
\end{enumerate}
\end{defbox}

\begin{propbox}[Proposition --- Relation Between Geometric and Algebraic Multiplicity]
\label{prop:mult-inequality}
Let $f$ be a linear operator on a finite-dimensional vector space $V$ over
$\F$ and let $\lambda$ be an eigenvalue of $f$.  Then
\[
    1 \;\le\; \dime V_\lambda \;\le\; \operatorname{mult}(\lambda,\chi_f(x)).
\]
\end{propbox}

% ---- Equivalent conditions for diagonalizability ------------
\begin{thmbox}[Theorem --- Equivalent Conditions for Diagonalizability]
\label{thm:diag-equiv}
Let $f:V\to V$ be a linear operator on a finite-dimensional vector space
$V$ over $\F$, and let $\lambda_1,\ldots,\lambda_k$ be the distinct
eigenvalues of $f$.  The following are equivalent:
\begin{enumerate}[label=\textbf{(\roman*)}]
    \item $f$ is diagonalizable.
    \item The characteristic polynomial splits as
          $\chi_f(x)=\prod_{i=1}^k(x-\lambda_i)^{d_i}$ and
          $\dime V_{\lambda_i}=d_i$ for each $i$.
    \item $\displaystyle\dime V = \sum_{i=1}^k \dime V_{\lambda_i}$.
\end{enumerate}
Here $d_i=\operatorname{mult}(\lambda_i,\chi_f(x))$ is the algebraic
multiplicity of $\lambda_i$.
\end{thmbox}

\begin{propbox}[Proposition --- $n$ Distinct Eigenvalues Implies Diagonalizable]
\label{prop:n-distinct-eigenvalues}
Let $V$ be an $n$-dimensional vector space over $\F$ and let $f$ be a
linear operator on $V$.  If $\chi_f(x)$ has $n$ distinct roots in $\F$,
then $f$ is diagonalizable.
\end{propbox}
